\section{Limitaciones}

\begin{enumerate}
  \item \textbf{No hubo ajuste fino end-to-end en la corrida nueva.} La máquina
  disponible no expuso CUDA. Se compararon backbones ImageNet congelados con
  una cabeza logística optimizada. Es un experimento completo sobre el corpus,
  pero no demuestra el rendimiento que tendrían las mismas arquitecturas al
  adaptar todos sus pesos.

  \item \textbf{La comparación con Etapa 1 no es controlada.} Cambiaron corpus,
  división, objetivo de entrenamiento y familia de cabeza. Las diferencias
  históricas ayudan a entender el proyecto; no aíslan el efecto causal de la
  ampliación o de Optuna.

  \item \textbf{Los grupos de entorno están confundidos con la clase.}
  Solamente roya, mancha gris y tizón norteño tienen imágenes tanto de
  laboratorio como de campo. El gap global resume el dataset, no mide por sí
  solo el efecto de retirar un fondo uniforme.

  \item \textbf{Faltan metadatos de campo.} No se conoce finca, municipio,
  variedad, teléfono, campaña o evaluador. La fuente del archivo no reemplaza
  esos datos y puede estar asociada simultáneamente con clase, resolución y
  estilo de captura.

  \item \textbf{Blur y resolución son indicadores indirectos.} La varianza del
  Laplaciano depende del contenido de la lesión y la compresión, no solamente
  del enfoque. Los terciles describen diferencias internas y no fijan un
  umbral universal de calidad.

  \item \textbf{La validación cruzada es condicional a una selección previa.}
  Arquitectura, hiperparámetros y ensemble se eligieron con la división
  train/validación. Los cinco folds posteriores miden estabilidad dentro del
  desarrollo; el estimador verdaderamente no consultado es el holdout final.

  \item \textbf{No se calibraron probabilidades.} La confianza sirve para
  ordenar alternativas y activar una advertencia de interfaz, pero no se
  comprobó confiabilidad mediante ECE, diagrama de calibración o temperatura
  ajustada en un conjunto separado.

  \item \textbf{La explicación del ensemble es aproximada.} Grad-CAM combina
  mapas normalizados de componentes usando los pesos de soft voting. Esto ayuda
  a localizar regiones influyentes, pero no produce una descomposición causal
  exacta de la probabilidad conjunta. LIME usa 32 perturbaciones en un caso y
  explica solo EfficientNet-B0, el componente de mayor peso.

  \item \textbf{La latencia no es móvil.} Los valores comparativos de backbones
  y la demo corresponden a CPU de escritorio. El artefacto móvil histórico sí
  está integrado, pero no se ejecutó un benchmark en varios teléfonos Android
  durante esta etapa.

  \item \textbf{Disponibilidad pública no confirmada al cierre.} La aplicación
  aprobó typecheck y 276 pruebas, pero el dominio de la API dejó de resolver en
  la comprobación final y la URL de APK respondió 404 sin autenticación. No se
  otorgan todos los puntos de despliegue.

  \item \textbf{Cobertura diagnóstica limitada.} Nueve etiquetas no representan
  todas las enfermedades, plagas, daños químicos, estrés hídrico ni síntomas
  mixtos. El sistema no determina dosis, causalidad ni manejo a partir de una
  sola foto.

  \item \textbf{No hay validación prospectiva salvadoreña ni medición de
  impacto.} Los resultados provienen de colecciones públicas. No se midió
  concordancia con especialistas en parcelas, ahorro, rendimiento, adopción o
  comprensión del usuario.
\end{enumerate}

Ninguna de estas limitaciones invalida los artefactos generados. Define con
precisión qué pregunta respondieron. La etapa demuestra un protocolo de
selección reproducible y una ruta funcional de producto; todavía no demuestra
eficacia agronómica ni disponibilidad operativa a escala.
